\section{Introduction}
\IEEEPARFLOWSTART{U}{rban} freight transportation and last-mile delivery operations are currently undergoing two simultaneous structural transformations: rapid fleet electrification and algorithmic digitization~\cite{toth2014vehicle}. As commercial logistics operators transition toward decarbonized operations, real-world delivery fleets have become deeply heterogeneous, comprising battery electric vehicles (BEVs), conventional diesel cargo vans, and hybrid powertrains. 

Crucially, in operational settings, vehicles cannot be treated as nominal identical entities. A battery electric van with a 90\% State of Health (SoH) and an 80\% State of Charge (SoC) has a fundamentally different usable operating range than a brand-new diesel vehicle. Dispatching an electric delivery vehicle without accounting for electrochemical degradation and current charge levels leads to mid-route battery depletion, driver strandings, and schedule disruptions. Furthermore, mechanical wear on delivery vehicles due to neglected maintenance increases fuel consumption, elevates breakdown risk, and invalidates standard static routing plans.

Simultaneously, the combinatorial optimization community has investigated hybrid metaheuristics~\cite{burke2013hyper,talbi2009metaheuristics} and quantum approximate optimization algorithms~\cite{farhi2014qaoa,hadfield2019qaoa,feld2019hybrid} to tackle the NP-hard Capacitated Vehicle Routing Problem (CVRP). However, our forensic audit of existing hybrid literature reveals three pervasive structural deficiencies:
\begin{enumerate}
    \item \textbf{Pseudo-Cooperation:} Most proposed "cooperative" engines merely run algorithms independently and compare their final outputs, or execute a rigid, unidirectional pipeline ($A \to B \to C$). There is no persistent shared candidate population, no cross-algorithm exchange, and no mechanism allowing one solver to iteratively improve a candidate generated by another paradigm~\cite{toulouse2004cooperative}.
    \item \textbf{Isolated Quantum Toy Problems:} Quantum routing literature (e.g., QAOA~\cite{azad2022solving} and quantum annealing~\cite{feld2019hybrid}) focuses almost exclusively on tiny symmetric graphs ($N \le 8$ nodes) with zero connection to enterprise Big Data pipelines or practical fleet management constraints.
    \item \textbf{Vehicle Condition Obliviousness:} Existing routing algorithms optimize pure spatial distance under rigid, uniform vehicle capacity assumptions, ignoring vehicle-specific powertrain costs, energy consumption, and maintenance degradation.
\end{enumerate}

To overcome these barriers, this paper presents \textbf{ROUTE/IQ-NextGen}, an adaptive, vehicle-condition-aware cooperative optimization framework. Rather than running solvers in isolation, ROUTE/IQ-NextGen orchestrates seven heterogeneous operators---Apache Spark polar sweeps, Google OR-Tools, simulated QAOA variational boundary cuts, Quantum-Inspired Genetic Algorithm (QIGA) rotation gates, Tabu Search, Deep Variable Neighborhood Descent (VND), and Guided Local Search---around a shared candidate pool blackboard.

The remainder of this paper is structured as follows. Section~\ref{sec:related_work} reviews the 27 foundational papers. Section~\ref{sec:problem_formulation} establishes the HFVRP-VC mathematical formulation. Section~\ref{sec:system_architecture} details the system architecture. Section~\ref{sec:big_data} covers Big Data processing and BPR congestion modeling. Section~\ref{sec:quantum_classical} describes the quantum and quantum-inspired operators. Section~\ref{sec:cooperative_engine} elaborates on the candidate exchange blackboard and Contextual MAB. Section~\ref{sec:vehicle_aware} details the vehicle condition evaluation. Section~\ref{sec:math_proof} gives the complete MILP model and the formal proof of Proposition~1. Section~\ref{sec:methodology} presents the experimental methodology. Section~\ref{sec:results} discusses empirical benchmark and ablation findings. Section~\ref{sec:contributions} states research contributions. Section~\ref{sec:limitations} provides honest scientific disclosures. Section~\ref{sec:conclusion} concludes the paper.
